\section{Basic concepts of Lie algebra}

Let \(L\) be a Lie algebra over a field \(F\).

\begin{definition}
    The \emph{derived series} of \(L\) is defined inductively by
    \[
        L^{(0)}=L, \qquad L^{(1)}=[L,L], \qquad L^{(i)}=[L^{(i-1)},L^{(i-1)}]\quad (i\ge 1).
    \]
    The Lie algebra \(L\) is called \emph{solvable} if
    \[
        L^{(n)}=0
    \]
    for some integer \(n\ge 0\).
\end{definition}

\begin{definition}
    The \emph{descending central series} of \(L\) is defined by
    \[
        L^1=L, \qquad L^2=[L,L], \qquad L^i=[L,L^{i-1}]\quad (i\ge 2).
    \]
    The Lie algebra \(L\) is called \emph{nilpotent} if
    \[
        L^n=0
    \]
    for some integer \(n\ge 1\).
\end{definition}

\begin{proposition}
    If \(L\) is nilpotent, then \(L\) is solvable.
\end{proposition}

\begin{proof}
    It is enough to show that
    \[
        L^{(i)}\subseteq L^{2^i}\qquad (i\ge 0).
    \]
    We argue by induction on \(i\). The case \(i=0\) is clear. If
    \[
        L^{(i)}\subseteq L^{2^i},
    \]
    then
    \[
        L^{(i+1)}
        =
        [L^{(i)},L^{(i)}]
        \subseteq
        [L^{2^i},L^{2^i}]
        \subseteq
        L^{2^{i+1}}.
    \]
    Hence the claim follows for all \(i\). If \(L^n=0\) for some \(n\), then for \(i\) sufficiently large one has \(2^i\ge n\), so
    \[
        L^{(i)}\subseteq L^{2^i}=0.
    \]
    Therefore \(L\) is solvable.
\end{proof}

\begin{definition}
    A \emph{solvable ideal} of \(L\) is an ideal which is solvable as a Lie algebra in its own right. Since the sum of two solvable ideals is again solvable, there exists a unique maximal solvable ideal of \(L\), called the \emph{radical} of \(L\), denoted by
    \[
        \operatorname{Rad}(L).
    \]
    The Lie algebra \(L\) is called \emph{semisimple} if
    \[
        \operatorname{Rad}(L)=0.
    \]
\end{definition}

\begin{remark}
    Thus the three notions fit into the chain
    \[
        \text{nilpotent} \;\Longrightarrow\; \text{solvable},
    \]
    while semisimplicity means that \(L\) has no nonzero solvable ideals. In particular, a semisimple Lie algebra cannot contain any nonzero abelian ideal.
\end{remark}

\begin{example}
    Any abelian Lie algebra is nilpotent, hence solvable. On the other hand, the Lie algebra of upper triangular matrices is solvable in general, but need not be nilpotent.
\end{example}

\begin{example}
    Let
    \[
        L=
        \left\{
        \begin{pmatrix}
            a & b \\
            0 & 0
        \end{pmatrix}
        \,\middle|\,
        a,b\in F
        \right\}
        \subseteq \mathfrak{gl}_2(F),
    \]
    Then \(L\) is solvable but not nilpotent.

    Indeed, if
    \[
        x=
        \begin{pmatrix}
            a & b \\
            0 & 0
        \end{pmatrix},
        \qquad
        y=
        \begin{pmatrix}
            c & d \\
            0 & 0
        \end{pmatrix},
    \]
    then
    \[
        [x,y]
        =
        \begin{pmatrix}
            0 & ad-bc \\
            0 & 0
        \end{pmatrix}.
    \]
    Hence
    \[
        [L,L]
        =
        \left\{
        \begin{pmatrix}
            0 & t \\
            0 & 0
        \end{pmatrix}
        \,\middle|\,
        t\in F
        \right\},
        \qquad
        [[L,L],[L,L]]=0.
    \]
    Thus the derived series terminates:
    \[
        L^{(1)}=[L,L],\qquad L^{(2)}=[L^{(1)},L^{(1)}]=0,
    \]
    so \(L\) is solvable.

    On the other hand, \(L\) is not nilpotent, since its descending central series does not reach \(0\). In fact,
    \[
        L^2=[L,L]
        =
        \left\{
        \begin{pmatrix}
            0 & t \\
            0 & 0
        \end{pmatrix}
        \,\middle|\,
        t\in F
        \right\},
    \]
    and for any
    \[
        x=
        \begin{pmatrix}
            a & b \\
            0 & 0
        \end{pmatrix}\in L,\qquad
        y=
        \begin{pmatrix}
            0 & t \\
            0 & 0
        \end{pmatrix}\in L^2,
    \]
    we have
    \[
        [x,y]
        =
        \begin{pmatrix}
            0 & at \\
            0 & 0
        \end{pmatrix}\in L^2.
    \]
    Therefore
    \[
        L^3=[L,L^2]=L^2,
    \]
    and inductively
    \[
        L^n=L^2\neq 0\qquad (n\ge 2).
    \]
    So \(L\) is not nilpotent.
\end{example}

